\chapter{Variational Dynamics}

\section{What Argmin Was Really Asking}

Chapter 14's projection operator was stated as a minimization: Y = \(\operatorname*{argmin}_{Y \in C}\) ‖Y − (\(M_t\) + \(\Delta\psi\))‖². At the time, this was introduced as a single-step computation, find the nearest admissible point to a given proposal, without asking whether it belonged to anything larger. It is worth asking now. A minimization performed once, at a single moment, in isolation from everything else the field does, is a coincidence if it is not secretly an instance of a more general principle governing the field's behavior throughout. This chapter argues it is not a coincidence, and supplies the principle.

\section{An Energy Functional for the Field}

Suppose the field's coherence, flow, and entropy jointly determine a single scalar quantity, an energy, at every configuration the field could take. Write this as a functional H, built from a Lagrangian density evaluated across the field's domain:

\[ L = \tfrac{1}{2}\|v\|^2 - \Phi - \delta S + \eta\langle v, \nabla\Phi\rangle. \]

Each term earns its place by referring back to a chapter already written. The kinetic-like term, ½‖v‖², costs energy in proportion to how much flow is occurring at a point, penalizing motion for its own sake unless something justifies it. The potential-like term, −Φ, rewards coherence: high-Φ regions, the settled, self-reinforcing configurations Chapter 16 called dynamic fixed points, sit at lower energy than fragmented ones, which is exactly the sense in which this book has always treated coherence as something the field tends toward rather than away from. The entropy term, −\(\delta S\), rewards resolution over ambiguity, consistent with Chapter 8's insistence that entropy is a real cost the field carries rather than a neutral bookkeeping quantity. The coupling term, η⟨v, ∇Φ⟩, rewards flow that moves toward increasing coherence and penalizes flow that moves away from it, which is the energetic expression of something Chapter 10 already described in different language: that attention and action follow the gradient of activated relevance, itself grounded in the field's own coherence structure. H, the field's total energy, is the integral of L across the whole domain, and it is worth being clear that nothing about this functional is new content. It is a single object that several chapters' worth of separately motivated tendencies turn out to be describing pieces of.

\section{Intrinsic Relaxation as Gradient Descent}

With H in hand, Chapter 13's \(F_{\mathrm{phys}}\) can finally be given more than a name. The field's intrinsic tendency to relax, to smooth coherence, to let entropy settle, to let flow subside once nothing continues to drive it, is exactly the negative functional gradient of H:

\[ F_{\mathrm{phys}}(M) = -\delta H/\delta M. \]

This is worth dwelling on, because it closes a claim Chapters 13 and 14 made without fully justifying it. Section 14.4 argued that \(F_{\mathrm{phys}}\) is close to admissible almost by construction, because it is built from the same discipline that defines C. It is now possible to say something stronger: \(F_{\mathrm{phys}}\) does not merely happen to respect C, it is literally the direction of steepest energy decrease, and to the extent C's hard constraints were chosen, in Chapter 18, to characterize exactly the configurations this energy functional favors, \(F_{\mathrm{phys}}\) moves the field toward lower energy and toward greater admissibility in the same motion, because for a well-constructed C these are close to the same thing. Gradient descent on H is not a mechanism this book is introducing for the first time in this chapter. It is what Chapter 13 was already calling \(F_{\mathrm{phys}}\), given its proper variational name.

\section{Projection as Constrained Minimization}

Chapter 14's argmin can now be recognized as a special case of the same underlying principle, restricted rather than unrestricted. Unconstrained gradient descent on H would move a proposal toward whatever configuration minimizes energy globally, ignoring C entirely. Constraint projection instead minimizes a different, more local functional, the squared distance to the proposal, subject to remaining inside C. These are not two unrelated computations that happen to share the word minimization. They are the same variational apparatus applied to two different objectives, one minimizing global energy without regard for admissibility, the other minimizing distance-from-proposal subject to admissibility, and a full treatment of constrained variational problems, of the kind this chapter will not develop in detail, shows that the second can itself be recovered as an unconstrained minimization of H plus a penalty term that grows without bound as a configuration moves outside C. Projection, in other words, is what gradient descent on H looks like once C's hard constraints are folded into the energy functional as an infinite wall rather than left as a separate condition checked afterward.

\section{Fixed Points as Critical Points}

This gives Chapter 16's fixed-point equation, and the stability question Chapter 16 explicitly deferred, their proper mathematical content. A fixed point \(M^*\) satisfying \(M^*\) = \(\Pi_C\)(\(M^*\) + \(\Delta\psi\)) is, in the variational picture, a critical point of H restricted to C: a configuration where the projected gradient vanishes, where there is no direction of further descent available without leaving the admissible region. A stable fixed point, in Chapter 16's informal sense, that small perturbations are absorbed back toward \(M^*\) rather than amplified away from it, is precisely a local minimum of H restricted to C, a configuration surrounded on all admissible sides by higher energy, so that any small displacement costs energy the field's own dynamics will spend the next several updates recovering. An unstable fixed point, Chapter 16's precarious stack of loose stones, is a critical point that is not a local minimum, a saddle or local maximum where some admissible direction actually decreases energy further, so that the smallest nudge in the right direction sends the field's dynamics away from \(M^*\) rather than back toward it. Chapter 16 promised this distinction would receive its proper formal footing in Part IV. This is that footing.

\section{What Guarantees Existence}

Chapter 18 closed by flagging a real gap: nothing established there guaranteed that a nearest point in C actually exists for a given proposal, only that the question was well posed once C was given a genuine geometric definition. The variational picture developed in this chapter does not close that gap outright, but it says precisely what closing it would require. The standard tools of the calculus of variations show that a minimizer exists whenever the functional being minimized is bounded below and does not admit a sequence of ever-decreasing values escaping to infinity without converging, conditions usually called coercivity, and whenever the admissible set is closed under the limits of such sequences. Applied here, this becomes a concrete, checkable claim about H and C rather than an appeal to intuition: if H is coercive on C and C is closed in the appropriate sense, then a nearest admissible point exists for every proposal, and Chapter 14's argmin was never at risk of returning nothing. This book will not carry out that verification in the main text; it belongs, together with the fuller apparatus of Lagrange multipliers this chapter has gestured at without deriving, to the extended proofs this book's appendices reserve for exactly this purpose.

\section{Looking Ahead}

This chapter has treated the field's dynamics as though they were purely dissipative: energy decreases, proposals settle, fixed points are where descent finally stops. This is not the whole picture, and it is worth being honest about what it leaves out before moving on. The fountain from Chapter 16 was offered as a dynamic fixed point precisely because its water never stops moving even while its form remains stable, and pure gradient descent, taken alone, describes a field settling toward rest, not a field sustaining perpetual, patterned motion without ever losing energy to it. Something in this book's account needs to explain how a region can cycle indefinitely, water through a fountain, flow through a courtyard's drainage, without that cycling being merely energy draining slowly toward zero. That requires a genuinely different kind of dynamics, one where energy is conserved rather than dissipated, and Chapter 20 is where this book introduces it.
